% HiMCM 2024 Problem A --- Task 6 letter to the IOC Executive Board
% Compile from this directory: pdflatex ioc_letter.tex
% Policy claims below follow results/brisbane_2032_ranking.csv and
% results/sensitivity_weight_shock.csv. Target length: one and a half pages.
\documentclass[12pt]{letter}

\usepackage[margin=1in]{geometry}
\usepackage{setspace}
\usepackage[hidelinks]{hyperref}

\hypersetup{
  pdftitle={Letter to the IOC Executive Board: Brisbane 2032 SDE shortlist},
  pdfauthor={GoHiMCM 2024 A study team}
}

\signature{Programme Evaluation Working Group\\GoHiMCM 2024 Problem~A study team}
\address{GoHiMCM Studies\\Olympic Programme Evaluation\\For the IOC Executive Board}
\date{19 September 2026}

\begin{document}
\begin{letter}{The IOC Executive Board\\
International Olympic Committee\\
Ch\^{a}teau de Vidy\\
1007 Lausanne\\
Switzerland}

\opening{Dear Members of the Executive Board,}

\setstretch{1.05}

\paragraph{Executive summary.}
We assessed three sports already in public discussion for Brisbane 2032---Flag
football, Cricket, and Squash---against the same six IOC tests the programme
always faces: popularity and access, event-load discipline, worldwide reach,
gender equity, relevance to young people, and the sustainability of venues.
On that scorecard, and using only the locked team priorities already applied
to historical sports, the Board should treat \textbf{Flag football as first
preference} (composite 0.536), \textbf{Cricket as second} (0.484), and
\textbf{Squash as third} (0.396). Twenty-percent swings in any one policy
weight leave this medal order unchanged. The recommendation is therefore
stable under ordinary disagreement, but it is \emph{not} a blank cheque:
Flag football's lead collapses if youth appeal is taken off the table, and
Cricket remains the heaviest venue proposition of the three.

\paragraph{Model credibility.}
Before we ranked 2032 candidates, the same framework was run on sports the
IOC has already added, kept, or set aside; it placed recent urban additions
(skateboarding, breaking, surfing) above karate's one-Games experiment, while
still scoring athletics and swimming among the strongest long-stay sports---in
line with the direction of programme decisions since Tokyo 2020.

\paragraph{Actionable recommendations for Brisbane 2032.}
\textit{First --- Flag football.} It is the only candidate that combines a
full gender split (50~percent) with the strongest youth mark in the trio and
a modest venue burden. For Olympic legacy it offers a light-footprint team
sport that can sit on existing rectangular fields, speak to a school-age
audience, and stay inside a two-event medal load. We advise a Brisbane
listing of one women's and one men's tournament, with no extra formats in
2032.

\textit{Second --- Cricket.} It is the broadcast leader of the shortlist and
the only candidate with a repeated Olympic appearance in the contest history
file, and it is the natural host-nation story in Australia. Its contribution
to legacy is continental reach and prime-time audiences, especially if the
Board keeps the programme to a small medal event and uses grounds Australia
already operates. It should not be sold as a low-cost add-on: among the three
it has the highest infrastructure score.

\textit{Third --- Squash.} It is the reach leader (150 federations) and the
cheapest venue profile, with a near-parity gender mark. It does less for a
youth-first narrative than Flag football, which is why it finishes third
under current priorities. Its Olympic contribution is universality and an
indoor, compact spectator product that can be staged as a temporary show
court rather than a permanent hall.

\paragraph{Risk and mitigation for 2036 and beyond.}
Two stresses appeared in the sensitivity file. First, \textbf{youth-preference
fragility:} when youth weight is set to zero, Squash rises to first and Flag
football falls to third; on a wider historical panel the same test lifts
athletics and swimming and drops breaking. Mitigation: write Flag football's
2032 entry as a youth-and-gender contract (equal events, school pathway,
strict event cap). If a future Board downgrades youth as a criterion, promote
Squash from ``third'' to the default indoor add, rather than expanding Flag
football's medal count. Second, \textbf{Cricket's venue cost:} do not approve
a purpose-built Olympic oval. Require a lean hosting protocol---existing
international grounds, overlay rather than new builds, and a single-city or
tight two-ground plan so 2036 hosts are not locked into stadium debt. Squash
should travel as a plaza or arena overlay; Flag football should remain a
field overlay. None of the three should be allowed to grow event numbers as
the price of remaining on the programme.

We recommend that the Board (i)~invite Flag football for Brisbane 2032 on
the two-event, overlay basis above; (ii)~keep Cricket in the 2032
conversation only with a written no-new-stadium condition; (iii)~hold Squash
as the robustness reserve for 2032 if youth priorities shift, and as the
preferred compact add for 2036+ if venue pressure rises.

\closing{Yours faithfully,}

\end{letter}
\end{document}
